Table~\ref{tab:torus-period-n60} shows the revival periods
observed on each torus, for $N \le 60$. A bullet marks each
$(k, m)$ and $N$ where at least one canonical revival exists.
The distribution is highly uneven. The tori involving $k = 2$
contribute the smallest families, because the $2$-cycle has
only three signatures, all at $N_{\min} \le 6$. The largest
families occur on the tori $4 \times 4$, $3 \times 4$, and
$3 \times 3$, where the dense 1D spectra of the $3$- and
$4$-cycles allow many parameter combinations to share a common
period. The period $N = 60$ appears on the most tori, because
the $5$- and $10$-cycles revive only at $60$ and any torus with
a $5$ or $10$ axis therefore inherits that period.
